% =============================================================================
% CONTROLE SINTÉTICO (16:37 a 16:49)   [Script 07: exemplo completo]
% =============================================================================
\section[Sintético]{Controle sintético}

\begin{frame}{A ideia: fabricar um controle}
  \centering
  \includegraphics[height=3.6cm]{07_sc_trajetorias.png}\par
  \footnotesize Um único estado adotou o programa em 2010 (hipotético). \textbf{Nenhum doador} (cinza) imita o Estado T (laranja).
  \vspace{0.1cm}
  \begin{itemize}
    \item Constrói-se um \ind{Estado T sintético}: média \emph{ponderada} de estados não tratados que reproduz a trajetória do Estado T \textbf{antes} da política;
    \item O impacto é a \destaque{distância} entre os dois \textbf{depois}.
  \end{itemize}
\end{frame}

\begin{frame}{Controle sintético: características e exemplos}
  \carac{combinação \textbf{ponderada} de unidades não tratadas que imita o tratado no pré}
        {bom \textbf{ajuste pré-política}; sem choques só do tratado depois; sem contaminação dos doadores}
        {efeito sobre a unidade tratada, \textbf{ano a ano}}
        {painel \textbf{longo} (muitos anos antes), poucas unidades tratadas, bom grupo de doadores}
        {leis e programas de um estado, país ou município}
  \vspace{0.05cm}
  \textbf{\footnotesize\color{azul}Exemplos em políticas públicas}\par
    \scriptsize
    \begin{itemize}
      \item \textbf{EUA:} lei antitabaco da Califórnia (Proposição 99), com os outros estados como doadores (Abadie, Diamond e Hainmueller, 2010).
      \item \textbf{Espanha:} custo econômico do terrorismo no País Basco (Abadie e Gardeazabal, 2003).
      \item \textbf{Hipotético:} programa estadual de prevenção à violência, resultado = homicídios por 100 mil (script 07, com o passo a passo completo).
    \end{itemize}
\end{frame}

\begin{frame}{O ajuste antes da política e a lacuna depois}
  \centering
  \includegraphics[height=4.2cm]{07_sc_ajuste.png}\par
  \footnotesize Antes de 2010, o sintético (39\% do Doador 16, 24\% do Doador 8\ldots) acompanha o Estado T. Depois, a distância média é de \destaque{$-4{,}6$} homicídios por 100 mil (verdade: $-5$).
\end{frame}

\begin{frame}{Placebos: a lacuna do Estado T se destaca?}
  \centering
  \includegraphics[height=4.2cm]{07_sc_placebos.png}\par
  \footnotesize O mesmo método aplicado a cada doador (cinza). Só o Estado T (laranja) se afasta depois de 2010: 1 em 20 unidades, $p=0{,}05$.
\end{frame}
